
\chapter{Combat}



\textit{Combat} refers to an ordered, turn-based conflict. 
This conflict is typically martial, 
but may also be social (e.g., negotiating with a guard, 
or haggling with a shopkeeper).



\section{Teams}
 Relationships between characters may be described as \textbf{Hostile}, 
 \textbf{Neutral}, or \textbf{Allied}.

\vspace{0.4em}

 Allied characters may join a battle as a \textit{Team}.
 When a team attacks another group of characters that are Allied with one another, 
 then the defending characters become a team, and the teams become Hostile to one another.



\section{Rounds \& Turns}
Combat \textit{Rounds} consist of teams taking \textit{Turns} 
to perform \textit{Actions}.


\subsection{Round Structure}
Within a round, teams take turns in the order they join the battle. 
The aggressor joins the battle first, 
followed by the team defending against the attack.
Other teams are given an option to join or flee after the defender 
has taken their turn.

\vspace{0.4em}

Teams take their turns simultaneously, freely ordering their actions.


\subsection{Turn Structure}
When it becomes a character's team's turn to act, 
it is considered the ``beginning" of their turn.

\vspace{0.4em}

If multiple effects would trigger at the beginning of a turn,
the team whose turn it is resolves their triggered effects first, 
in any order they choose, followed by the remaining teams in turn order.

\vspace{0.4em}

Each round, when it is a character's turn, they may make 3 \textit{Actions}.

\vspace{0.4em}

At any point in a round, a character may use a \textit{Reaction} ability
if the conditions for that ability are met, 
but they may only use one Reaction each round.

\vspace{0.4em}

Once a character uses an ability as an action or reaction, 
it cannot be used again until the beginning of their next turn.




\section{Effects}
\textit{Effects} are rules that apply a change to the game state, or occasionally, to the rules themselves.

\vspace{0.4em}

When an effect would contradict the core rules, the rule introduced by the effect takes precedence.
When an effect would contradict another effect, the effect that was activated most recently takes precedence.

\subsection{Activated Effects}
\textit{Activated} effects are effects that may be applied by the character at their discretion,
within the limitations described by the effect.

\vspace{0.4em}

These effects are generally applied once using an \textit{Action}, and after the effect duration expires
(which may be instantaneous or may span a number of actions or turns),
may be reapplied or renewed by reactivating the effect with another action.

\vspace{0.4em}

Activated effects that are available for use by characters are generally referred to as ``Abilities''.

\subsection{Passive Effects}
\textit{Passive} effects are effects that apply continuously to or from the source of the effect,
depending on the specific rule, and are only removed while negated, 
or when the source of the effect is itself destroyed or otherwise removed.

\vspace{0.4em}

Passive effects that apply to characters are generally referred to as ``Perks''.

\subsection{Triggered Effects}
\textit{Triggered} effects may be active or passive, 
and are only evaluated when some condition is met by the game state.

An activated effect with a trigger gives the character the option to use the ability 
when the condition is met, whereas a passive effect with a trigger is automatically applied
when the condition is met.


\newpage

\section{Core Abilities}
These abilities are available to all characters and most creatures, with a few exceptions.


\subsection{Movement}
When a character uses the \textit{Move} action, 
they must designate a \textit{Movement Path} that is up to the length 
of their movement distance.

\vspace{0.4em}

If the path would intersect within base contact to a Hostile character, 
that character may use their Reflex reaction (see below) to intercept.
If the intercepting character deals Damage to the moving character, 
the movement must end within base contact.

\vspace{0.4em}

When the path of movement passes through or ends in hazardous terrain, 
roll on the Table (\Roll{d10}) for the terrain effect.

\vspace{0.4em}

Characters can move vertically (climb) at half-speed, 
but ending movement on a vertical surface will return the Character 
to the foot of the terrain.


\subsection{Equip/Unequip}
If a character has enough free hands available, 
they may \textit{Equip} into their hand(s) an item that is within base contact 
or in their inventory.

\vspace{0.4em}

If a character has an item equipped in their hand, 
they may \textit{Unequip} that item and return it to an empty slot in their inventory
or place it within base contact.


\subsection{Use Item}
A character may use an ability of an equipped item, such as a weapon or potion.


\subsection{Interact}
A character may trigger the ability or effect of an non-equippable object
that is within base contact.


\subsection{Prepare}
Choose an action ability your character has not used this turn. 
Until their next turn, if they do not use that ability as an action, 
they may use that ability as a reaction instead.


\subsection{Reflex}
When a Hostile character ends their movement in base contact, 
or leaves base contact by moving away,
your character may use their Reaction to make a normal action 
as though they had used Prepare for that ability.


\newpage

\subsection{Attack}
Weapons and some abilities can \textit{Attack} an enemy target as part of their effect.
To make an Attack:

\begin{enumerate}
  \item Select an eligible target in range of the attack ability.
  \item Perform a Clash Roll against the target.
  \item If victorious in the Clash, roll to deal Damage.
\end{enumerate}


\subsection{Grapple}
To \textit{Grapple}, a character must be in base contact 
with the defending character and have an open hand.

\vspace{0.4em}

If the defending character is \textbf{Similar-size or Smaller}:
\begin{enumerate}
  \item The characters Clash, with the aggressor adding their 
  Strength modifier, 
  and the defender adding either their Agility or Strength modifier.
  \item If the aggressor wins:
    \begin{itemize}
      \item Aggressor may equip the defender as an unwieldy 
      (unusable) 2-handed item
      \item If the characters are Similarly-sized, 
      Aggressor's move distance becomes 1"
    \end{itemize}
\end{enumerate}

If the defending character is \textbf{Larger}:
\begin{enumerate}
  \item The characters Clash, 
  with the aggressor adding either their Agility or Strength modifier 
  and losing Advantage, 
  and the defender adding either their Agility or Strength modifier.
  \item If the aggressor wins:
    \begin{itemize}
      \item Aggressor may mount the defender as an unruly 
      (uncontrollable) steed.
    \end{itemize}
\end{enumerate}

If a character wins a grapple while equipped or mounted by the 
defending character, they become unequipped or unmounted by the defender.




\newpage

\section{Tactics \& Positioning}

\subsection{Formation}
Similarly-sized characters on the same team and 
contiguously within 1" are considered to be in a \textit{Group}.

\vspace{0.4em}

If all characters in a group would make identical attacks
against another group within common range, 
you may instead have them attack and defend together.
(Identical attacks use the same rolls and modifiers, 
and trigger the same effects.)

To attack as a group:
\begin{enumerate}
  \item Combine the attacks into a single Clash roll that is rolled as a 
    \Roll{d6} pool.
  \item Defending group rolls the same number of \Roll{d6}.
  \item Pair the dice pools in descending order.
  \item Resolve the corresponding pairs as normal Clash rolls, 
    with each side adding the lowest Clash modifier of all characters 
    in the group.
  \item Defender assigns unsaved attacks as they choose 
    among the members of the group.
  \item Attacker rolls Damage for each unsaved attack. 
    If a character is killed before all damage rolls for unsaved attacks 
    have been resolved, Defender reassigns the remaining unsaved attacks, 
    and Attacker continues rolling Damage.
\end{enumerate}



\subsection{High Ground}
Two characters at different elevations of sheer (vertical) 
terrain are considered in base contact if:
\begin{itemize}
  \item the characters' bases are in contact with the ledge and 
    foot of the terrain.
  \item the separation between bases is less than the width of the 
    base of the character at the foot of the terrain.
\end{itemize}

If the characters are Hostile, 
the character in contact with the ledge gains the \textit{High Ground},
while the character in contact with the foot gains the \textit{Low Ground}.

\vspace{0.4em}

Attacks made with the High Ground against the target 
with the Low Ground gain Advantage on all rolls.

\vspace{0.4em}

Attacks made from the Low Ground against the target 
with the High Ground lose Advantage on all rolls.



\subsection{Flanking}
When an attack deals damage to a target, 
all characters that are within range to attack the target,
and can draw a straight line from their base, through the enemy model, 
to the base of the ally that was source of the damage, 
gain \textit{Flanking} against the target until the start of their next turn.

\vspace{0.4em}

While a character is Flanking a target, 
they gain Advantage on Clash rolls for attacks against that target.



\subsection{Cover}
Projectile attacks with partially-obscured Line-of-Sight grant \textit{Cover} 
to the target for the duration of the attack.

\vspace{0.4em}

Attacks against enemies with cover lose Advantage on Clash Rolls.



\section{Damage Types \& Resistances}
Some sources of damage can support a \textit{Damage Type} that adds additional effects
beyond the amount of damage is dealt.

\subsection{Physical Damage}
If a melee or projectile weapon could deal more than one type of \textit{Physical} damage,
choose a damage type to use when making the attack:

\begin{itemize}
  \item Slashing
  \item Piercing
  \item Chopping
  \item Bludgeoning
\end{itemize}

\subsection{Elemental Damage}
If a weapon attack could deal \textit{Elemental} damage,
that effect occurs as a separate roll after any effect triggered by the physical damage type.

\vspace{0.4em}

Elemental effects of all types that are applied to an attack occur, 
and trigger in the order chosen by the source of the damaging attack.

\begin{Table}{c X}
  \TableHeader{Type} & \TableHeader{Effect} \\
  Fire & Deal \Roll{1d4} additional damage. On a Critical Damage roll for this element, the target is inflicted with the Ignite condition for \Roll{1d10} turns. \\
  Spark & Deal \Roll{1d4} damage to the target's Mana. On a Critical Damage roll for this element, the target is inflicted with the Electrocuted condition for \Roll{1d10} turns. \\
  Frost & Roll \Roll{1d10} for a cumulative effect: On a 1--5, there is no additional effect; on a 6+, the target's maximum movement is reduced by 1" for \Roll{1d10} turns (to a minimum of 1"); on a 9+, the target's maximum actions-per-turn is reduced by 1 for the duration of the effect (to a minimum of 1); on a 10+, the target is afflicted with the Frozen condition for the duration of the effect. \\
\end{Table}


\subsection{Damage Resistances}
Characters with resistance to a damage type cannot be affected by any additional effects 
inherent to that damage type,
and cannot be affected by any effects that are triggered by that damage type.

\vspace{0.4em}

Damage rolls against the resistant character that use that type of damage lose Advantage.



\section{Conditions}
\textit{Conditions} are passive effects that are applied to a character by another effect,
and may impact the character's abilities or effectiveness in the game.

\subsection{Prone}
If a character's movement distance is greater than 1", 
they may choose to become Prone as part of a Move action.

\vspace{0.4em}

If a character is Prone, their maximum movement distance becomes 1",
and they may use a Move action to remove the Prone condition.

\vspace{0.4em}

A character that is Prone receives Cover against projectiles,
and Hostile characters in Base contact with the character gain High Ground against them,
while the Prone character gains the Low Ground against the Hostile characters.

\subsection{Unconscious}
The character cannot take turns, benefit from passive effects, or travel.

\subsection{Paralyzed}
The character cannot use abilities, and cannot travel.

\subsection{Dazed}
The character's \textbf{Wisdom}, \textbf{Intelligence}, and \textbf{Zest} are reduced by 5,
to a minimum of 0.

\subsection{Restrained}
The character cannot use abilities requiring free hands or items equipped in the hands.
The character's \textbf{Agility}, \textbf{Dexterity}, and \textbf{Strength} are reduced by 5,
to a minimum of 0.

\subsection{Blinded}
The character cannot make Story rolls requiring sight.
The character cannot use abilities that target characters beyond the range of Base.
Abilities with a range of Base lose Advantage twice.

\subsection{Deafened}
The character cannot make Story rolls requiring hearing.
The character loses Advantage on all rolls against all characters that are Flanking them.

\subsection{Bleeding}
The character takes \Roll{1d4} damage for every action they make.

\vspace{0.4em}

This condition is removed when the character's Life is restored by any amount greater than 1.

\subsection{Crippled}
The character receives a passive effect depending on the body part that is Crippled:

\begin{Table}{c X}
  \TableHeader{Body Part} & \TableHeader{Effect} \\
  Right/Left Leg & The character's maximum movement distance is reduced to 1", and they may not use more than one action per turn to perform movement. The character's \textbf{Agility} is reduced to 0. \\
  Right/Left Arm & The corresponding hand becomes unequipped, and they cannot use abilities requiring the use of that free hand. The character's \textbf{Dexterity} is reduced to 0.  \\
  Torso & The character's maximum Life is reduced to 10, and their \textbf{Strength} and \textbf{Resilience} are reduced to 0. \\
  Head & The character is Unconscious for the duration of this effect, and their \textbf{Wisdom}, \textbf{Intelligence}, and \textbf{Zest} are reduced to 0. \\
\end{Table}

If a body part that is Crippled becomes inflicted with Crippled again, 
this effect becomes permanent and can only be removed when they level-up in a Town with a Hospital.
Otherwise, this effect is removed when the character is restored to their maximum Life.

\subsection{Exhausted}
A character may be inflicted with Exhaustion any number of times.
The character's maximum actions-per-turn is reduced by 1, to a minimum of 0. 
When a character is inflicted with Exhaustion, 
if it would reduce their maximum actions-per-turn below 0, instead they become Unconscious.

\subsection{Ignited}
The character takes \Roll{1d4} damage at the start of their turn.

\vspace{0.4em}

This condition is removed when the character takes Frost damage.

\vspace{0.4em}

If the character's maximum movement distance is greater than 1" and they are Prone, 
they may take a Move action to remove this condition.

\subsection{Frozen}
The character is Paralyzed for the duration of this effect.

\vspace{0.4em}

This condition is removed when the character takes Fire damage.

\subsection{Electrocuted}
All of the character's abilities gain an additional Mana cost of 1.




